% !TeX root = ../MAIN.tex
% !TeX encoding = UTF-8

\chapter*{Kurzfassung}

Die Erstellung von Dienstplänen für Pflegekräfte ist eine komplexe Aufgabe, die organisatorische Anforderungen und individuelle Präferenzen der Mitarbeitenden berücksichtigen muss. Diese Bachelorarbeit beschreibt die Entwicklung eines ASP-basierten Systems zur automatisierten Erstellung und kurzfristigen Anpassung von Dienstplänen in der Augenabteilung eines Krankenhauses, umgesetzt als vollständige Fullstack-Webanwendung mit Angular, FastAPI und PostgreSQL.
Die Optimierungslogik wird durch den ASP-Solver Clingo realisiert, der harte Constraints wie Mindestbesetzung und Qualifikationsanforderungen sowie weiche Constraints in Form einer fairen Stationsverteilung modelliert. Zusätzlich unterstützt das System ein Rescheduling-Verfahren für kurzfristige Ausfälle sowie die vorausschauende Berechnung von Alternativplänen. Die Arbeit zeigt, dass ASP eine geeignete Grundlage für praxistaugliche Planungssoftware im Gesundheitssektor darstellt.